\subsection{Обратимость возмущенного оператора \(A+\Delta A\)}

\begin{theorem}[Обратимость возмущенного оператора]
	\textbf{Условия:}
	\begin{itemize}
		\item $E$ --- банахово пространство;
		\item $A \in \mathcal{L}(E)$ --- обратимый оператор, причем
		$A^{-1} \in \mathcal{L}(E)$;
		\item $\Delta A \in \mathcal{L}(E)$;
		\item выполнено условие
		\[
		\|A^{-1}\|\,\|\Delta A\| < 1.
		\]
	\end{itemize}
	
	\textbf{Следовательно:}
	\begin{itemize}
		\item оператор $A+\Delta A$ обратим;
		\item $(A+\Delta A)^{-1} \in \mathcal{L}(E)$.
	\end{itemize}
\end{theorem}

\begin{proof}
	\begin{enumerate}
		\item \textbf{Сведение к оператору вида $I+B$.}
		Положим
		\[
		B := A^{-1}\Delta A \in \mathcal{L}(E).
		\]
		Тогда
		\[
		\|B\|
		\leqslant
		\|A^{-1}\|\,\|\Delta A\|
		<1,
		\]
		и
		\[
		A+\Delta A
		=
		A(I+A^{-1}\Delta A)
		=
		A(I+B).
		\]
		
		\item \textbf{Ряд Неймана.}
		Рассмотрим операторный ряд
		\[
		S := \sum_{n=0}^{\infty}(-B)^n.
		\]
		Так как
		\[
		\|B^n\|\leqslant\|B\|^n
		\]
		и числовой ряд
		\[
		\sum_{n=0}^{\infty}\|B\|^n
		\]
		сходится, то ряд $\sum\limits_{n=0}^{\infty}(-B)^n$
		сходится в банаховом пространстве $\mathcal{L}(E)$.
		Следовательно, $S\in\mathcal{L}(E)$.
		
		\item \textbf{Обратимость $I+B$.}
		Для частичной суммы
		\[
		S_N=\sum_{n=0}^{N}(-B)^n
		\]
		имеем
		\[
		(I+B)S_N=S_N(I+B)=I+(-1)^N B^{N+1}.
		\]
		Поскольку
		\[
		\|B^{N+1}\|\leqslant\|B\|^{N+1}\to0,
		\]
		при $N\to\infty$ получаем
		\[
		(I+B)S=S(I+B)=I.
		\]
		Значит,
		\[
		(I+B)^{-1}=S\in\mathcal{L}(E).
		\]
		
		\item \textbf{Обратимость $A+\Delta A$.}
		Так как $A$ и $I+B$ обратимы,
		\[
		(A+\Delta A)^{-1}
		=
		(I+B)^{-1}A^{-1}.
		\]
		Причем оба множителя принадлежат $\mathcal{L}(E)$, поэтому
		\[
		(A+\Delta A)^{-1}\in\mathcal{L}(E).
		\]
	\end{enumerate}
\end{proof}

\begin{corollary}
	В условиях теоремы
	\[
	(A+\Delta A)^{-1}
	=
	\sum_{n=0}^{\infty}
	(-A^{-1}\Delta A)^n A^{-1}.
	\]
	Кроме того,
	\[
	\|(A+\Delta A)^{-1}\|
	\leqslant
	\frac{\|A^{-1}\|}
	{1-\|A^{-1}\|\,\|\Delta A\|}.
	\]
\end{corollary}

\begin{proof}
	Из доказательства теоремы
	\[
	(I+B)^{-1}
	=
	\sum_{n=0}^{\infty}(-B)^n,
	\qquad
	B=A^{-1}\Delta A.
	\]
	Поэтому
	\[
	(A+\Delta A)^{-1}
	=
	(I+B)^{-1}A^{-1}
	=
	\sum_{n=0}^{\infty}
	(-A^{-1}\Delta A)^n A^{-1}.
	\]
	
	Также
	\[
	\|(I+B)^{-1}\|
	\leqslant
	\sum_{n=0}^{\infty}\|B\|^n
	=
	\frac{1}{1-\|B\|}.
	\]
	Следовательно,
	\[
	\|(A+\Delta A)^{-1}\|
	\leqslant
	\frac{\|A^{-1}\|}
	{1-\|A^{-1}\|\,\|\Delta A\|}.
	\]
\end{proof}

\begin{idea}
	\textbf{Суть:} Обратимость устойчива относительно достаточно малых
	возмущений. Если $A$ обратим и
	\[
	\|\Delta A\| < \frac{1}{\|A^{-1}\|},
	\]
	то $A+\Delta A$ также имеет ограниченный обратный оператор.
\end{idea}